\documentclass[11pt]{article}

\usepackage[T1]{fontenc}
\usepackage[utf8]{inputenc}
\usepackage{amsmath}
\usepackage{amssymb}
\usepackage{amsthm}
\usepackage{booktabs}
\usepackage{geometry}
\geometry{a4paper, margin=25mm}
\usepackage{hyperref}

\newtheorem{remark}{Remark}

\newcommand{\E}{\mathbb{E}}
\newcommand{\Var}{\operatorname{Var}}
\newcommand{\cv}{\mathrm{cv}}
\newcommand{\sd}{\mathrm{sd}}
\newcommand{\MSE}{\mathrm{MSE}}
\newcommand{\se}{\mathrm{se}}

\title{Two constants that were being dropped\\
\large the multiplicative constant in \texttt{prop:opt}'s allocation,
and the quantile in our confidence intervals}
\author{\texttt{loglog\_validation} project notes}
\date{2026-08-24}

\begin{document}
\maketitle

\begin{abstract}
Two independent findings from Experiments B and C, both of the same shape: a
constant that a theorem legitimately discards, and that costs real accuracy
when the theorem is applied literally. Part~\ref{part:kappa} derives, in closed
form, the multiplicative constant $\kappa$ that Proposition \texttt{prop:opt}
drops from its budget allocation --- worth a factor $8$--$38$ in compute ---
and shows it is recoverable from measured quantities alone, with no appeal to
ground truth. Part~\ref{part:cov} explains why every ``$95\%$'' confidence
interval this project published was in fact an $88\%$ interval, and derives the
figure $0.8784$ that the calibration experiment then measured as
$0.877$--$0.878$.
\end{abstract}

\tableofcontents

\part{The allocation constant \texorpdfstring{$\kappa$}{kappa}}
\label{part:kappa}

\section{What is being computed, and from what}

Proposition \texttt{prop:opt} (eq.~945--946) states that, given a compute
budget $B$, the choice
\begin{equation}
  n = \kappa B^{\theta_1}, \qquad m_0 = \theta_2\log_\rho B,
  \qquad \theta_1 = \frac{2\omega_1}{d+2\omega_1},
  \qquad \theta_2 = \frac{1}{d+2\omega_1},
  \label{eq:propopt}
\end{equation}
achieves the error decay $\lvert\hat\gamma-\gamma\rvert \lesssim
B^{-\omega_1/(d+2\omega_1)}$. It is a \emph{rate} result: correct up to
constants, and the constant multiplying $B$ inside the logarithm is exactly
what it discards. Empirically that omission costs a factor $2.0$--$3.4$ in
RMSE at every budget from $10^4$ to $10^9$; since $\mathrm{RMSE}\sim
B^{-1/3}$, a factor $c$ in RMSE is a factor $c^3$ in budget, i.e.\ $8$--$38$
times the compute.

The strategy below is to write the three quantities that matter --- bias,
standard deviation, and cost --- as exact power laws in $(n,m_0)$ with
explicit constants, and then minimise. Only the constants require
measurement; the exponents are the paper's.

\begin{remark}[No ground truth is used]
The inputs are $a_1$ and $\omega_1$ (Experiment B), $d$ (the affine cost fit),
and $\cv$ (read off the samples). Every one is a measured quantity. The true
values of $\gamma$, $a_1$, $\omega_1$ for simple random walk are known, but
they enter only when \emph{scoring} a finished answer, never when computing
one.
\end{remark}

\section{Bias}

The article's estimator is $\hat\beta=\sum_k w_{k,m}\log\overline Y_{\rho^k}$
with $\hat\gamma=\hat\beta/\log\rho$, on the consecutive grid
$k=m_0+1,\dots,m_0+m$, weights (eq.~526)
\begin{equation}
  w_{k,m}=\frac{12\bigl(k-m_0-\tfrac{m+1}{2}\bigr)}{m(m^2-1)} .
\end{equation}
Substituting the one-correction truncation of eq.~(232),
\begin{equation}
  \log\E Y_{\rho^k}=\log a_0+\gamma\, k\log\rho+a_1\rho^{-k\omega_1},
\end{equation}
and using the weight identities of eq.~(542), $\sum_k w_{k,m}=0$ and
$\sum_k w_{k,m}\,k=1$, the first two terms vanish identically and only the
correction survives:
\begin{equation}
  \E\hat\gamma-\gamma=\frac{a_1}{\log\rho}\sum_k w_{k,m}\rho^{-k\omega_1}.
\end{equation}
Now write $k=m_0+j$ with $j=1,\dots,m$. The weights depend on $j$ alone,
\begin{equation}
  w_j=\frac{12\bigl(j-\tfrac{m+1}{2}\bigr)}{m(m^2-1)},
\end{equation}
and $\rho^{-k\omega_1}=\rho^{-m_0\omega_1}\rho^{-j\omega_1}$ factors, so the
dependence on $m_0$ separates completely:
\begin{equation}
  \boxed{\;
  \lvert\text{bias}\rvert=C_b\,\rho^{-\omega_1 m_0},
  \qquad
  C_b=\frac{\bigl\lvert a_1\sum_{j=1}^{m}w_j\rho^{-j\omega_1}\bigr\rvert}{\log\rho}.
  \;}
  \label{eq:Cb}
\end{equation}
This is the finite-size bias of Prop.~(820), with its constant made explicit.

\section{Standard deviation}

By the delta method, $\Var(\log\overline Y_k)\approx
\sigma_k^2/(n\mu_k^2)=\cv^2/n$ whenever the coefficient of variation
$\cv=\sd(Y_i)/\E(Y_i)$ is scale-free (it is for $\lvert S_k\rvert$: measured
$0.7575$--$0.8222$ across the grid, against the half-normal limit
$\sqrt{\pi/2-1}=0.7555$). Distinct scales use independent samples, so
\begin{equation}
  \Var(\hat\gamma)=\frac{1}{\log^2\rho}\sum_j w_j^2\,\frac{\cv^2}{n}
  =\frac{\cv^2\lVert w\rVert^2}{n\log^2\rho},
  \qquad
  \lVert w\rVert^2=\sum_{j=1}^m w_j^2=\frac{12}{m(m^2-1)},
\end{equation}
the last equality from $\sum_j\bigl(j-\tfrac{m+1}{2}\bigr)^2=m(m^2-1)/12$.
Hence
\begin{equation}
  \boxed{\;
  \sd=C_s\,n^{-1/2},
  \qquad
  C_s=\frac{\cv\,\lVert w\rVert}{\log\rho}.
  \;}
  \label{eq:Cs}
\end{equation}

\begin{remark}[Consistency with the article's CLT]
Since $m(m^2-1)\to m^3$, this reproduces eq.~(583)'s
$\Var(\hat\gamma)=12\sigma_\infty^2/(nm^3\log^2\rho)$ --- an independent check
that $C_s$ carries the right normalisation.
\end{remark}

\section{Cost}

Under Assumption \texttt{cost\_is\_power\_law} (eq.~353), $\mathrm{cost}(i)=i^d$,
Lemma \texttt{lem:budget} gives the geometric sum in closed form:
\begin{equation}
  B=n\sum_{j=1}^{m}\rho^{d(m_0+j)}=n\,\rho^{dm_0}G,
  \qquad
  G=\rho^{d}\,\frac{\rho^{dm}-1}{\rho^{d}-1}.
  \label{eq:budget}
\end{equation}

\section{The optimisation}

Eliminating $n$ via \eqref{eq:budget}, i.e.\ $n=B\rho^{-dm_0}/G$, the mean
squared error becomes a function of $m_0$ alone:
\begin{equation}
  \MSE(m_0)=\underbrace{C_b^2\rho^{-2\omega_1 m_0}}_{\text{falls in }m_0}
  \;+\;\underbrace{\frac{C_s^2G}{B}\,\rho^{dm_0}}_{\text{rises in }m_0}.
  \label{eq:mse}
\end{equation}
The entire tradeoff is visible here: a deeper ladder reduces finite-size bias
but costs more per replicate, hence buys fewer samples. Setting
$u=m_0\log\rho$ and differentiating,
\begin{equation}
  \frac{\mathrm d}{\mathrm du}\Bigl[C_b^2e^{-2\omega_1u}
  +\tfrac{C_s^2G}{B}e^{du}\Bigr]=0
  \;\Longleftrightarrow\;
  2\omega_1C_b^2e^{-2\omega_1u}=d\,\frac{C_s^2G}{B}e^{du},
\end{equation}
so that $e^{(d+2\omega_1)u}=\dfrac{2\omega_1C_b^2}{d\,C_s^2G}B$, and therefore
\begin{equation}
  \boxed{\;
  m_0^{\star}=\theta_2\log_\rho(\kappa B)
  =\underbrace{\theta_2\log_\rho B}_{\text{\texttt{prop:opt}}}
  \;+\;\underbrace{\theta_2\log_\rho\kappa}_{\text{the dropped offset}},
  \qquad
  \kappa=\frac{2\omega_1C_b^2}{d\,C_s^2\,G}.
  \;}
  \label{eq:kappa}
\end{equation}
The second term does not depend on $B$. That is precisely why a rate argument
may discard it --- and precisely why the gap between \texttt{prop:opt}'s $m_0$
and the empirical optimum is a \emph{constant} rather than a growing
discrepancy.

Substituting $m_0^\star$ back into \eqref{eq:budget}, and using
$1-d\theta_2=\theta_1$,
\begin{equation}
  n^{\star}=\frac{B^{\theta_1}}{G\,\kappa^{d\theta_2}},
\end{equation}
which also identifies the unnamed constant $\kappa$ appearing in
\eqref{eq:propopt}.

\subsection{Two structural corollaries}

\paragraph{The optimal bias-to-noise ratio.}
The stationarity condition $2\omega_1\,\text{bias}^2=d\,\sd^2$ rearranges to
\begin{equation}
  \frac{\lvert\text{bias}\rvert}{\sd}\Big|_{m_0^\star}=\sqrt{\frac{d}{2\omega_1}}
  \;=\;0.7071 \quad\text{for }(d,\omega_1)=(1,1),
  \label{eq:ratio}
\end{equation}
\emph{independently of $B$}. This is the sharpest available diagnostic. Measured
at the empirical argmin it reads $0.44$--$1.50$ (the scatter is grid
granularity: one step of $m_0$ moves the ratio by $\rho^{\omega_1+d/2}=2.83$).
Measured at \texttt{prop:opt}'s own $m_0$ it reads $0.05$--$0.19$ --- an order
of magnitude below the balance its own derivation argues for. \texttt{prop:opt}
drives the bias far beneath the noise floor and pays for it in samples.

\paragraph{Why the constant is forgiving.}
Since $\kappa\propto C_b^2\propto a_1^2$ and $\kappa\propto C_s^{-2}\propto
\cv^{-2}$, and $\kappa$ enters \eqref{eq:kappa} only through a logarithm,
\begin{equation}
  \Delta(\text{offset})=2\theta_2\log_\rho\Bigl\lvert\frac{a_1'}{a_1}\Bigr\rvert
  =-2\theta_2\log_\rho\frac{\cv'}{\cv}.
\end{equation}
At $(d,\omega_1,\rho)=(1,1,2)$ this is $\pm 2/3$ of a step per doubling. Being
one step off costs at most $1.19\times$ in RMSE, so a factor-two error in $a_1$
is nearly free. The exception is $\omega_1$, which enters $\theta_2$, $C_b$ and
$\kappa$ simultaneously: halving it moves the offset by $2.12$ steps. It is
also the noisiest quantity we measure, which is the reason Experiment B pools
its replicates rather than averaging their fits.

\section{The numbers}

With $(d,\omega_1,\rho,m)=(1,1,2,6)$ and the measured $a_1=-0.2748$,
$\cv=0.7771$:
\begin{gather}
  w_j=(-0.14286,\,-0.08571,\,-0.02857,\,0.02857,\,0.08571,\,0.14286),\\
  \textstyle\sum_j w_j=2.8\times10^{-17},\qquad
  \sum_j j\,w_j=1.000000,\qquad
  \lVert w\rVert^2=0.057143=\tfrac{12}{6\cdot 35},\\
  S:=\textstyle\sum_j w_j\rho^{-j\omega_1}=-0.089732,
\end{gather}
\begin{align}
  C_b&=\frac{\lvert-0.2748\times(-0.089732)\rvert}{0.693147}=0.035575,\\
  C_s&=\frac{0.7771\times0.239046}{0.693147}=0.267999,\\
  G  &=2\cdot\frac{2^{6}-1}{2-1}=126,
\end{align}
\begin{equation}
  \kappa=\frac{2(1)(0.035575)^2}{(1)(0.267999)^2(126)}=2.7969\times10^{-4},
  \qquad
  \text{offset}=\tfrac13\log_2\bigl(2.7969\times10^{-4}\bigr)=-3.9346 .
\end{equation}

\section{Verification against Experiment C}

Sweeping $m_0$ at each budget ($R=40$ replicates per cell, so that every other
$m_0$ acts as a control arm for the one \texttt{prop:opt} names):

\begin{center}
\begin{tabular}{rrrrrr}
\toprule
$B$ & $m_0$ \texttt{prop:opt} & $m_0$ tuned & $m_0$ measured & penalty \texttt{prop:opt} & penalty tuned\\
\midrule
$10^4$ & 4 & 0 & 0 & $3.35\times$ & $1.00\times$\\
$10^5$ & 5 & 2 & 1 & $2.59\times$ & $1.01\times$\\
$10^6$ & 6 & 3 & 2 & $2.29\times$ & $1.00\times$\\
$10^7$ & 7 & 4 & 4 & $2.02\times$ & $1.00\times$\\
$10^8$ & 8 & 5 & 4 & $2.69\times$ & $1.02\times$\\
$10^9$ & 9 & 6 & 6 & $2.26\times$ & $1.00\times$\\
\bottomrule
\end{tabular}
\end{center}

\noindent
Penalty is RMSE relative to the best $m_0$ available on the grid. Three
observations:
\begin{enumerate}
  \item The gap $m_0^{\texttt{prop:opt}}-m_0^{\text{measured}}$ regressed on
    $\log_\rho B$ has slope $-0.05$ over five decades: it is flat, confirming
    that $\theta_2$ is right and only the constant was missing.
  \item The closed-form penalty for being $\delta$ steps too deep,
    $\sqrt{(r^2\rho^{-2\omega_1\delta}+\rho^{d\delta})/(r^2+1)}$ with
    $r^2=d/2\omega_1$, predicts $2.31\times$ at $\delta=3$ and $3.27\times$ at
    $\delta=4$; measured $2.02,2.26$ and $3.35,2.59,2.29,2.69$.
  \item After tuning, the residual penalty is $\le 2\%$ even where $m_0$ is off
    by one, because $\MSE$ is flat near its minimum: at $B=10^6$ the two best
    cells differ by $0.03\%$.
\end{enumerate}

The RATE half of \texttt{prop:opt} is unaffected and passes: the measured decay
exponent is $-0.3259\pm0.0116$ at the best $m_0$ against a predicted
$-0.3307\pm0.0127$, i.e.\ agreement to $0.28$ standard errors.

\part{Why our ``\texorpdfstring{$95\%$}{95\%}'' intervals were
\texorpdfstring{$88\%$}{88\%} intervals}
\label{part:cov}

\section{The interval as it was computed}

Experiment B reports a parameter $\theta\in\{\omega_1,a_1,\gamma\}$ as
\begin{equation}
  \hat\theta \;\pm\; 1.96\,\widehat{\se},
  \qquad
  \widehat{\se}=\frac{s}{\sqrt R},
  \qquad
  s^2=\frac{1}{R-1}\sum_{r=1}^{R}\bigl(\hat\theta_r-\bar\theta\bigr)^2,
\end{equation}
where $\hat\theta_1,\dots,\hat\theta_R$ are the fits from $R=5$ independent
replicates. (The point estimate $\hat\theta$ is the \emph{pooled} refit rather
than $\bar\theta$; see Remark~\ref{rem:pooled}.)

\section{The error}

The multiplier $1.96$ is the standard normal quantile, appropriate when
$\widehat\se$ is the \emph{true} standard error. Here it is not: $s$ is
estimated from five numbers. If $\hat\theta_r\sim\mathcal
N(\theta,\sigma^2)$ i.i.d., the exact pivot is Student's,
\begin{equation}
  T=\frac{\bar\theta-\theta}{s/\sqrt R}\;\sim\;t_{R-1}=t_4,
  \qquad\text{not }\mathcal N(0,1),
\end{equation}
so the true coverage of the nominally-$95\%$ interval is
\begin{equation}
  \boxed{\;
  \Pr\bigl(\lvert t_4\rvert<1.9600\bigr)=2F_{t_4}(1.9600)-1=0.8784 .
  \;}
\end{equation}

Two distinct effects push in the same direction. First, $s$ is
\emph{downward-biased}: $\E[s]=c_4(R)\,\sigma$ with
\begin{equation}
  c_4(R)=\sqrt{\frac{2}{R-1}}\;\frac{\Gamma(R/2)}{\Gamma\bigl((R-1)/2\bigr)},
  \qquad c_4(5)=0.9400,
\end{equation}
so a five-point standard deviation reads $6\%$ low on average. Second, $s$ is
\emph{noisy}, which fattens the tails of $T$ beyond the normal. The $t_{R-1}$
quantile corrects both simultaneously.

\section{Correct multipliers}

\begin{center}
\begin{tabular}{lrrr}
\toprule
nominal level & normal & $t_4$ & ratio\\
\midrule
$95\%$   & $1.9600$ & $2.7764$ & $1.417$\\
$68.3\%$ & $1.0000$ & $1.1417$ & $1.142$\\
\bottomrule
\end{tabular}
\end{center}

\noindent
The widely-used ``$\hat\theta\pm 1\,\se$'' shorthand is a $68.3\%$ claim, and
suffers the same defect:
$\Pr(\lvert t_4\rvert<1)=2F_{t_4}(1)-1=0.6261$.

\section{Measured coverage}

\texttt{src/check\_coverage.py} replays Experiment B's exact configuration
(scales $8..256$, the real per-scale $n$, $R=5$) $2000$ times against the known
ground truth for $\lvert S_k\rvert$ and counts interval hits. Interval bounds on
the coverage itself are Wilson score intervals.

\begin{center}
\begin{tabular}{llcc}
\toprule
quantity & centre & $q=$ normal & $q=t_4$\\
\midrule
$\omega_1$ & pooled       & $0.880$ $[0.866,0.894]$ & $0.948$ $[0.937,0.956]$\\
$\omega_1$ & mean-of-fits & $0.877$                 & $0.946$\\
$a_1$      & pooled       & $0.880$                 & $0.950$\\
$a_1$      & mean-of-fits & $0.878$                 & $0.955$\\
$\gamma$   & pooled       & $0.891$                 & $0.952$\\
$\gamma$   & mean-of-fits & $0.878$                 & $0.950$\\
\bottomrule
\end{tabular}
\end{center}

\noindent
The mean-of-fits rows match the predicted $0.8784$ to three decimal places, as
they must: for that centre the pivot really is $t_4$. At the $68.3\%$ level the
same agreement holds, $0.625$ measured for $\omega_1$ against $0.6261$
predicted.

\subsection{Two deviations, both explicable}

\begin{remark}[The pooled centre is mildly conservative]
\label{rem:pooled}
The pooled rows read slightly \emph{above} $0.8784$. The pooled estimate is not
$\bar\theta$, so its sampling distribution is not exactly $t_4$; pooling the
$\overline Y_i$ before fitting reduces the centre's scatter (for $a_1$:
$0.0349$ against mean-of-fits' $0.0390$) while the stated $\widehat\se$ is
unchanged, so the ratio $\widehat\se/\sd$ improves from $0.870$ to $0.974$ and
the interval becomes slightly wide rather than slightly narrow.
\end{remark}

\begin{remark}[Bias costs more in a narrow interval]
At the $68.3\%$ level, $a_1$/mean-of-fits reads $0.589$ against $0.6261$
predicted. That centre carries a residual bias of $-0.0155$ (the nonlinearity
of the fit, which pooling removes: pooled bias is $-0.0019$). A fixed bias
consumes a larger share of a $\pm1\,\se$ interval than of a $\pm1.96\,\se$ one.
\end{remark}

\section{Diagnosis, and what was changed}

The defect is the \emph{quantile}, not the width. The diagnostic
$\widehat\se/\sd(\hat\theta)$ measures $0.94$--$1.03$ for the pooled estimates:
the error bar is the right size. Nothing about the model, the estimator, or the
pooling was at fault.

Accordingly the printed $\pm$ values were left alone --- a standard error is a
perfectly good thing to report --- and the fix was applied where a standard
error becomes a \emph{decision}. The consistency test in
\texttt{src/plot\_allocation.py} used $\lvert z\rvert<2$, valid only for an
exactly known $\se$; here $z$ combines an analytic slope error with one derived
from $\omega_1$'s replicate error. Components are now combined by
Welch--Satterthwaite,
\begin{equation}
  \se^2=\sum_j\se_j^2,
  \qquad
  \nu_{\mathrm{eff}}=\frac{\bigl(\sum_j\se_j^2\bigr)^2}{\sum_j\se_j^4/\nu_j},
\end{equation}
and the cut-off is $t_{\nu_{\mathrm{eff}}}$. For Experiment C,
$\nu_{\mathrm{eff}}=16.9$ and the threshold is $2.111$ rather than $1.960$.
Both existing verdicts are unchanged, which is the desirable outcome: the
correction widens the threshold without overturning a conclusion.

\section{Practical consequence}

At $R$ replicates, use $t_{R-1}$. Note that $R=5$ is an expensive operating
point:
\begin{center}
\begin{tabular}{lrrrr}
\toprule
$R$ & 5 & 10 & 20 & $\infty$\\
\midrule
$t_{R-1}$ at $95\%$ & $2.776$ & $2.262$ & $2.093$ & $1.960$\\
excess width & $+42\%$ & $+15\%$ & $+7\%$ & --- \\
\bottomrule
\end{tabular}
\end{center}
\noindent
If replicates are cheap relative to the ladder itself, buying more of them is a
better use of budget than it appears.

Finally, this calibration is \emph{not} portable. Sweeping the sample size
shows coverage holding to $n$-scale $\approx 0.1$ and then collapsing --- not
because the interval degrades but because the estimator does: $(a_1,\omega_1)$
ceases to be identified once noise swamps the correction, and at $n$-scale
$0.003$ the fit returns $\hat a_1\approx-2\times10^{10}$. Coverage must be
re-measured before an error bar is trusted at a smaller budget, or on a
different model.

\appendix
\section{Reproducing the numbers}

\begin{verbatim}
# Part I: constants, tuned allocation, and the m0 sweep
python3 src/allocation_experiment.py \
    -meta experiments/01_srw/allocation_wide_config.json --tag allocation_wide
python3 src/plot_allocation.py -data experiments/01_srw/data/allocation_wide
python3 src/allocation_table.py --compare

# Part II: coverage calibration
python3 src/check_coverage.py --arm planted  --trials 2000 --centre both
python3 src/check_coverage.py --arm planting --trials 3000 --planting-n 10000
\end{verbatim}

\noindent
The constants themselves are \texttt{tools/allocation.py}'s
\texttt{allocation\_constants} / \texttt{tuned\_allocation}; the calibration
machinery is \texttt{tools/coverage.py}. Narrative write-ups with the full
context live in \texttt{experiments/01\_srw/README.md}.

\end{document}
